% ==========================================================
% 页面版式设置（纸张大小、页边距、页眉页脚间距）
% ==========================================================
\usepackage[
  a4paper,           % A4 纸张
  margin=2.5cm,      % 四周页边距
  headheight=14pt,   % 页眉高度
  headsep=10pt,      % 页眉与正文距离
  footskip=35pt      % 页脚高度
]{geometry}

% ==========================================================
% 表格相关宏包
% ==========================================================
\usepackage{longtable} % 跨页长表格
\usepackage{array}     % 增强列格式
\usepackage{booktabs}  % 三线表

% ==========================================================
% PDF 页面插入
% ==========================================================
\usepackage{pdfpages}  % 插入整页 PDF

% ==========================================================
% 正文字体大小与行距（推荐稳定版 + Markdown 列表兼容）
% ==========================================================
\fontsize{12pt}{22pt}\selectfont % 正文字号 12pt + 行距 22pt

\usepackage{setspace}  % 用于 Markdown 列表行距控制
% \setstretch{1.2}       % 正文 + 列表行距约 22~26pt
\setstretch{1.5}

% ==========================================================
% 中文字体设置
% ==========================================================
\usepackage{xeCJK}
\setCJKmainfont{SimSun}[BoldFont=SimHei, ItalicFont=KaiTi]    % 主字体
\setCJKsansfont{Microsoft~YaHei}[BoldFont=*~Bold]             % 无衬线
\setCJKmonofont{FangSong}                                      % 等宽字体

% ==========================================================
% 英文字体设置
% ==========================================================
\usepackage{fontspec}
\setmainfont{Times New Roman}  % 英文正文字体
\setmonofont{Consolas}         % 英文等宽字体（代码）

% ==========================================================
% 标题格式设置
% ==========================================================
\usepackage{titlesec}

% 四级标题 (\paragraph)
\titleformat{\paragraph}{\normalfont\normalsize\bfseries}{\theparagraph}{1em}{}
\titlespacing*{\paragraph}{0pt}{1ex plus 0.5ex}{1ex}

% 一级标题 (\section)
\titleformat{\section}{\fontsize{16pt}{22pt}\selectfont\bfseries}{\thesection}{1em}{}
\titlespacing{\section}{0pt}{13pt}{13pt}

% 二级标题 (\subsection)，后三个参数分别是左缩进，标题前垂直间距，标题后垂直间距
\titleformat{\subsection}{\fontsize{12pt}{18pt}\selectfont\bfseries}{\thesubsection}{1em}{}
\titlespacing{\subsection}{0pt}{10pt}{10pt}

% 三级标题 (\subsubsection)
\titleformat{\subsubsection}{\fontsize{12pt}{16pt}\selectfont\bfseries}{\thesubsubsection}{1em}{}
\titlespacing{\subsubsection}{0pt}{8pt}{4pt}

% ==========================================================
% 文档主标题 (\maketitle) 大小设置
% ==========================================================
% \makeatletter  % 允许使用 @ 命令
% \renewcommand{\maketitle}{%
%   \begin{center}%
%     {\Huge\bfseries \@title\par}%
%     \vspace{1em}%
%     \ifx\@author\@empty\else
%       {\Large \@author\par}%
%     \fi
%     \ifx\@date\@empty\else
%       \vspace{0.5em}%
%       {\large \@date\par}%
%     \fi
%   \end{center}%
%   \vspace{2em}%
%   \par\noindent% 确保退出垂直模式并开始正常段落
% }
% \makeatother   % 恢复 @ 的普通含义

% ==========================================================
% 目录与 PDF 书签
% ==========================================================
\usepackage[
  bookmarks=true,
  bookmarksopen=true,  % 展开书签
  bookmarksdepth=3,
  hidelinks,           % 隐藏链接边框
  % 以下似乎没用
  % colorlinks=true,     % 使用颜色标记链接，而不是边框（覆盖hidelinks）
  % linkcolor=blue,      % 内部引用颜色
  % citecolor=green,     % 引用文献颜色
  % urlcolor=red,        % URL链接颜色
  % linktoc=all,         % 目录也应用颜色，目录中的页码和标题都可点击（默认只标题可点）
]{hyperref}

% ==========================================================
% 图片与浮动体设置
% ==========================================================
\usepackage{graphicx}
\usepackage{float}
\usepackage{caption}

\setlength{\intextsep}{18pt}        % 文中浮动体上下间距
\setlength{\textfloatsep}{18pt}     % 页顶/页底浮动体与正文间距
\setlength{\abovecaptionskip}{8pt}  % caption 上方间距
\setlength{\belowcaptionskip}{8pt}  % caption 下方间距
\captionsetup{skip=8pt}             % caption 与图像间距

% --------------------------
% 封面页不显示页码
% --------------------------
\usepackage{etoolbox}
\AtBeginDocument{\thispagestyle{empty}}

% ==========================================================
% 页眉页脚设置
% ==========================================================
\usepackage{fancyhdr}
\pagestyle{fancy}
\fancyhf{}
% \fancyhead[L]{}                 % 左页眉
% \fancyhead[R]{}     % 右页眉
\fancyfoot[C]{第~\thepage~页}        % 页脚居中

% ==========================================================
% 列表层级与符号设置
% ==========================================================
\usepackage{enumitem}
\setlistdepth{4}                         % 支持 4 层列表
\setlist[itemize,1]{label=\textbullet, itemsep=6pt, topsep=6pt} 
\setlist[itemize,2]{label=\textopenbullet, itemsep=4pt, topsep=4pt} 
\setlist[itemize,3]{label=\(\triangleright\), itemsep=2pt, topsep=2pt} 
\setlist[enumerate,1]{itemsep=6pt, topsep=6pt}
\setlist[enumerate,2]{itemsep=4pt, topsep=4pt}
\setlist[enumerate,3]{itemsep=2pt, topsep=2pt}

% ==========================================================
% 代码块显示设置（listings）
% ==========================================================
% \usepackage{listings}
% \usepackage{xcolor}

% \definecolor{codebg}{RGB}{240,240,240}  % 浅灰色背景

% \lstset{
%   backgroundcolor=\color{codebg},
%   basicstyle=\ttfamily\small,
%   breaklines=true,
%   frame=single,
%   frameround=tttt,
%   numbers=left,
%   numberstyle=\tiny\color{gray},
%   keywordstyle=\color{blue},
%   commentstyle=\color{green!60!black},
%   stringstyle=\color{red!60!black},
%   showstringspaces=false,
%   xleftmargin=\parindent,  % 可以适当调整左侧边距
%   linewidth=\textwidth  % 设置代码块宽度不超过页面宽度
% }

\usepackage{fvextra}
\DefineVerbatimEnvironment{Highlighting}{Verbatim}{breaklines, breakanywhere=true, commandchars=\\\{\}}
\DefineVerbatimEnvironment{verbatim}{Verbatim}{breaklines, breakanywhere=true, commandchars=\\\{\}}